\textbf{Conciliación vigente INC-39.} La supervisión independiente añade cuatro EOLR-1 en extremos ordinarios, cuatro monitores y cuatro EOL a la selección desarrollada aquí: 56 monitores y 128 direcciones, cuatro SLC 40/40/24/24 con 88/88/52/52 mA. No cambia seis VEP, cuatro FRPS, ocho baterías 33 Ah ni cuatro 75 Ah. El controlador pasa a 1,232686/5,776 A y 36,078957 Ah de cribado; cada FRPS con relé a 0,625/0,751 A a 24 V y 18,0751 Ah. Los balances de INC-36 que siguen son la base antes de esos cuatro ramales, sin segunda lista activa. La selección y sus límites de subtensión, autonomía y supervisión están en INC/EXT-39.

\textbf{Detección láser y sectorización del sitio (INC-4).}\label{sec:sd4-inc4}
Synaptix selecciona seis detectores Xtralis VESDA-E VEP-A00-1P-UL: cuatro ordinarios F01/F02/F03/F06 y dos para bancos F04/F05 sujetos al expediente de aptitud de INC-30; se conservan tres repuestos fríos, dos ordinarios y uno de bancos. INC-36 sustituye los cuatro VLF-500-00-UL de la selección anterior por VEP con máximos publicados de 0,57 A en reposo y 0,59 A en alarma a 24 V. La ficha VLF y sus 410 mA nominales/490 mA de alarma quedan como antecedente de la sustitución, no como consumo máximo de la oferta. Los objetivos de sensibilidad y transporte se comprueban por cada red en ASPIRE para el modelo vigente; no se heredan límites de red o certificaciones de otra variante \parencites[RT-06.16--19, pp.~15--16]{pucv-transversal}[pp.~1--2]{l18-vlf500ul}[p.~2]{lote30-vep36106}.

\begin{longtable}{|p{0.09\linewidth}|p{0.26\linewidth}|p{0.18\linewidth}|p{0.31\linewidth}|}
\hline \textbf{Zona} & \textbf{Recinto independiente} & \textbf{Área / volumen bruto} & \textbf{Red de aspiración de proyecto} \\\hline
\endfirsthead
\hline \textbf{Zona} & \textbf{Recinto independiente} & \textbf{Área / volumen bruto} & \textbf{Red de aspiración de proyecto} \\\hline
\endhead
F01 & TI, comunicaciones y climatización común & 27,5 m$^2$ / 82,5 m$^3$ & Una rama de hasta 18 m, 12 puntos: ocho superiores, dos en retorno de gabinetes y dos en retorno climático. \\\hline
F02/F03 & UPS A / UPS B & 5,5 m$^2$ / 16,5 m$^3$ cada una & Una rama de hasta 8 m y cuatro puntos por recinto, incluidos retorno UPS y techo. \\\hline
F04/F05 & Bancos A / bancos B & 12 m$^2$ / 36 m$^3$ cada uno & Una rama de hasta 14 m y ocho puntos por recinto: cuatro superiores y cuatro sobre gabinetes, sin tomar aire exterior. \\\hline
F06 & Circulación frontal & 7,5 m$^2$ / 22,5 m$^3$ & Una rama de hasta 10 m y cuatro puntos superiores. \\\hline
\end{longtable}

Los seis volúmenes suman $82,5+2(16,5)+2(36)+22,5=210$ m$^3$; las redes reservan 72 m de tubería y 40 puntos. La circulación se separa mediante puertas y cerramientos para que sea una sexta zona efectiva; si el plano ejecutivo no materializa ese cierre, se integra su volumen con el recinto comunicado y se recalculan agente, boquillas y retención. No se declara simultáneamente una circulación abierta y una zona sellada. Se utiliza tubería de aspiración de 25 mm compatible con la ficha, red independiente por detector, sin mezclar aire entre bancos y TI. Los detectores F01/F02/F03/F06 permanecen en su zona, en posición accesible, fuera de captación inmediata de extracción, con escape a ese mismo volumen. Para F04/F05 se selecciona VEP-A00-1P-UL según INC-30, cuya ficha declara Class I Division 2 A/B/C/D, sin equiparar ese alcance con Zone 1 ni aprobar el circuito completo. Llevar el detector fuera del banco manteniendo una tubería que aspira H$_2$ no elimina el riesgo dentro del detector o del escape. Clasificación de área, certificado íntegro y condiciones de detector, retorno, confirmador y accesorios siguen pendientes. INC-36 conserva la selección condicionada y recalcula también las cuatro redes ordinarias; RT-06.16 permanece no completado para el sitio íntegro. Su ambiente se mantiene entre 0 y 38~$^\circ$C, sin condensación. Las penetraciones se sellan sin obturar la tubería; las posiciones se cotejan con gabinetes, insuflación y extracción en máxima renovación.

Se fijan umbrales iniciales propios Alert 0,05; Action 0,10; Fire 1 0,20 y Fire 2 0,30\,\% obs/m, con retardos respectivos de 10, 5, 0 y 0 s. Son ajustes de proyecto dentro del rango publicado, no sensibilidad garantizada por punto de muestreo. El diseño ASPIRE y la recepción comprueban transporte desde el punto más remoto en no más de 45 s, sensibilidad efectiva por orificio, caudales normales y obstrucción/rotura de tubería, con climatización y extracción en sus extremos. Action inicia prealarma local y remota; Fire 1 entra a la lógica de confirmación. El mapa vigente utiliza tres estados por aspirante (Action, Fire 1 y Fault) hacia monitores individuales; la disponibilidad de siete relés programables VEP no añade señales o direcciones sin asignación. Alert y Fire 2 conservan registro e indicación conforme a su configuración. AutoLearn se emplea con supervisión y registro de los ajustes; no puede elevar los umbrales de la oferta sin una nueva aceptación. Synaptix entrega memoria de red, plano de puntos, configuración firmada, registro de ensayo con humo de prueba autorizado y mantenimiento de filtros/aspirador; la actuación no depende de Internet. Los ajustes publicados VLF son antecedente; la ficha VEP y ASPIRE gobiernan la selección vigente \parencites[p.~2]{l18-vlf500ul}[pp.~1--2]{lote30-vep36106}.

\textbf{Panel, confirmación y alimentación autónoma (PCI-4).}\label{sec:sd4-pci4}
Se selecciona un Fike Cheetah Xi 10-068-R-2-L, versión 240 V/50--60 Hz, con fuente suplementaria 10-2474-2. La capacidad inicial de dos lazos se amplía mediante un SLM 10-2473 a cuatro lazos de 254 dispositivos y 100 mA por lazo. La central conserva control por zonas, actuadores de extinción, módulos de monitor/relé y sensores; no se confunde con el panel convencional de un solo riesgo. Las zonas F01/F02/F03/F06 reciben un sensor fotoeléctrico/térmico Fike 63-1059 con base aisladora 63-1060, independiente del VESDA. Se reservan dos sensores independientes adicionales para F04/F05, pero no se autoriza instalar esos modelos ordinarios en bancos sin acreditar clasificación de gas y compatibilidad; la misma condición alcanza módulos, contactos, actuadores y accesorios que entren en el área clasificada. Esta reserva de seis sensores en total no acredita selección apta para las dos zonas de bancos. Los dieciocho contactos VESDA, tres por zona, ingresan mediante dieciocho módulos Fike 55-046, con supervisión de circuitos y final de línea conforme al manual. Un Fire 1 VESDA y el sensor independiente de la misma zona deben coincidir para iniciar extinción automática; dos niveles del mismo VESDA no son dos detecciones independientes. Se instalan seis módulos de liberación Fike 55-053 y seis estaciones manuales de agente 20-1343. Cada zona incluye estación de aborto mantenido compatible con el sistema, señalización de predescarga y descarga, presostato de confirmación y supervisión de cilindro. INC-36 selecciona seis avisadores P2RL-SP y uno de reserva, doce kits DFIA e interruptores 02-12534 según EXT-30; aborto, cobertura/ajustes de avisadores y ensamblaje/listado final conservan verificaciones antes de adquisición \parencite[pp.~1--6]{l18-cheetah}.

Los lazos A y B separan físicamente las rutas y sus aislamientos: F02/F04 van por el lateral A, F03/F05 por B; F01/F06 conservan asignación individual sin compartir el disparo. No se programa una descarga simultánea de las seis zonas por una alarma de una sola. Se reservan, como mínimo, seis sensores, dieciocho monitores VESDA, seis liberadores, seis estaciones manuales, seis entradas de aborto, doce entradas de presión/supervisión de cilindro, doce entradas de compuertas, seis relés de actuación y cuarenta y ocho relés de monitor auxiliar: 124 direcciones previstas al incluir cuatro fallas FRPS de INC-36, con 52 monitores 55-046 y 54 relés 55-048. INC-36 desarrolla cuatro lazos: F02/F04 y F03/F05 con 40 direcciones/88 mA cada uno; F01 y F06 más dos fallas FRPS por lazo con 22/48 mA cada uno. La distribución anterior 61/61 cabía por direcciones pero excedía el límite de corriente SLC; queda sustituida. Los dos sensores aptos de bancos siguen pendientes y su presupuesto no es modelo seleccionado. En particular, el módulo Ethernet Fike 10-2627 comunica historia entre paneles supervisores Fike; su ficha no acredita SNMP/Modbus genérico ni su uso como anunciación primaria del edificio. La integración auxiliar descrita más adelante utiliza contactos separados \parencites[pp.~3--6]{l18-cheetah}[pp.~1--3]{l18-fike-ethernet}.

La protección local se respalda según INC-36 con cuatro baterías OEM de 12 V/75 Ah para controlador/SPS, en dos conjuntos de serie y dos gabinetes, más ocho de 12 V/33 Ah en cuatro pares y cuatro gabinetes de fuentes remotas. Cuatro FRPS alimentan un VEP ordinario cada una; los dos VEP de bancos quedan en Cheetah, uno por fuente con tres DFIA. El manual 06-356 revisión 11 limita cada fuente Cheetah a 1,5 A de reposo y 6 A de alarma, conjunto 3/12 A. INC-36 desarrolla controlador 1,230746/5,768 A y SPS 0,7954/3,45 A, frente al subtotal precedente INC-33; las capacidades de cribado 36,0223/23,2525 Ah no acreditan curva útil, recarga, tensión final ni ruta. Gabinetes, temperatura, acceso y evacuación conservan verificación, sin atribuir autonomía por etiqueta de batería \parencites{lote33-cheetah11}{lote33-frps}.

El ingreso único de PCI-4 y su monitor auxiliar se conmuta entre alimentación protegida A/B, preferentemente en L2, sin sumar dos veces la carga. La reserva de adquisición anterior de INC-4, pendiente de recalcular con cuatro FRPS y cargadores de INC-36, es 0,300 kW y 0,350 kVA en servicio normal, y 0,650 kW/0,750 kVA durante alarma o recarga, incluidos los cuatro lectores auxiliares. Como presupuesto superior se usan 12 A$\times$24 V = 288 W CC, 2,6 A$\times$24 V = 62,4 W de recarga y eficiencia exigida no menor de 0,80; se añaden 140 VA para lectores, tomando conservadoramente 140 W: $(288+62,4)/0,80+140=578$ W, por debajo de 650 W para el alcance anterior, sin incluir como demostrada la nueva alimentación FRPS. No es balance vigente completo. La ficha no declara ese consumo AC total ni factor de potencia: son límites de adquisición que deben verificarse con las fuentes cargando, régimen de alarma y dispositivos definitivos. La pérdida de UPS o del monitor auxiliar no elimina la detección, evacuación ni liberación locales durante la autonomía prescrita.

\textbf{Extinción por agente limpio y seguridad de bancos (EXT-4).}\label{sec:sd4-ext4}
Se selecciona la familia Fike FM-200/HFC-227ea, con válvula Impulse, seis conjuntos independientes de cilindro/tubería/boquillas, uno por zona. El agente es no conductor y no deja residuos; la ficha oficial anuncia aprobaciones UL/ULC/FM y distingue sus límites de aplicación y exposición. Eso no acredita automáticamente el listado UL del conjunto concreto, ni convierte la certificación ISO del fabricante en aprobación de instalación. Synaptix exige expediente UL del sistema completo con cilindros, válvulas, boquillas, actuadores y control compatibles; el instalador competente diseña conforme a NFPA 2001 y a NFPA 72 para detección/control, en las ediciones aplicables al proyecto aceptadas por la autoridad. No se sustituye inundación del recinto por un cartucho de rack \parencites[pp.~1--2]{l18-fm200}[RT-06.17, p.~15]{pucv-transversal}.

Se toman los volúmenes brutos de la tabla, sin deducir gabinetes ni inventar doble piso/plenum. Como antecedente de dimensionamiento, sustituido en cantidades por los llenados de EXT-30, se fijó concentración de proyecto de 8,0\,\% y temperatura mínima de 10~$^\circ$C. El manual Fike 06-439, revisión 3, sección 2, p.~2 (página 51 del archivo), determina $s=0,1269+0,0005T$ m$^3$/kg y $m=(V/s)C/(100-C)$, con $C$ expresada como porcentaje. A 10~$^\circ$C, $s=0,1319$ m$^3$/kg: resultan 54,3890 kg para F01; 10,8778 kg para cada UPS; 23,7334 kg para cada banco; y 14,8334 kg para F06, total preliminar 138,4448 kg. La oferta vigente selecciona 140 kg por cohorte, 280 kg con reemplazos, conforme a EXT-30. Manteniendo aquella masa preliminar, a 38~$^\circ$C el porcentaje aproximado es 8,7746\,\%. El cálculo sustituye la reserva por gas ideal y emplea la ecuación aproximada del fabricante, pero el manual consultado corresponde a recipientes EN y aprobación LPCB: su fórmula no acredita un ensamblaje UL. El 8,0\,\% sigue siendo valor de proyecto sujeto al combustible, concentración extintora y factores normativos; no se demuestra extinción de hidrógeno ni de refrigerante R32, y deben verificarse altitud/presión y exposición. Se reserva un conjunto de recarga/reemplazo local por zona, separado de los cilindros instalados, y se contrata reposición inmediata tras descarga \parencite[sec.~2, p.~2]{l21-fike-manual}.

La memoria hidráulica del fabricante fijará modelo/capacidad de cilindro, presión de almacenamiento, masa efectiva, diámetros/longitudes/accesorios, reparto, orificios y número de boquillas; demostrará descarga y concentración exigidas, inicialmente no más de diez segundos de descarga, retención inicial de diez minutos o el tiempo mayor exigido por el riesgo, y alivio de presión calculado según resistencia real de cerramientos. El ensayo de integridad del recinto comprueba fugas y tiempo de retención; una masa calculada no demuestra hermeticidad. Los cilindros se instalan en gabinete técnico exterior protegido, anclado, fuera de la ruta de evacuación, con seis circuitos independientes y reserva de frente de servicio de 1,20 m. Su huella y ruta de tuberías se incorporarán al plano ejecutivo al seleccionar los recipientes; no están acreditadas por el plano actual. Las válvulas de cierre/compuertas y el alivio no descargarán hacia circulación ocupada ni comunicarán gases entre A/B y TI.

\begin{longtable}{|p{0.20\linewidth}|p{0.66\linewidth}|}
\hline \textbf{Estado local} & \textbf{Actuación prescrita, independiente de WAN} \\\hline
\endfirsthead
\hline \textbf{Estado local} & \textbf{Actuación prescrita, independiente de WAN} \\\hline
\endhead
Action / prealarma & Aviso de prealarma en panel/RDU2 y monitor; la evacuación general emplea los NAC y P2RL-SP de INC-36, sin atribuir dos tonos automáticos al mismo equipo. Se mantiene refrigeración y extracción de bancos; no se corta una rama sana ni se dispara agente por prealarma única. \\\hline
Confirmación de fuego & Coincidencia VESDA Fire 1 más sensor independiente de la zona, o estación manual autorizada: evacuación, bloqueo de ingreso, salida libre mecánica y señalización. Se inicia retardo propio de 30 s, con aborto mantenido supervisado y registro. No se presupone que toda intervención humana quede completada en ese tiempo. \\\hline
F01, F02/F03, F06 & Se detiene recirculación afectada y se cierran entradas/salidas de aire de la zona, comprobando fin de carrera antes de liberación. Se aísla la fuente eléctrica del riesgo afectado mediante elemento diseñado para ello; se preservan panel, sirenas, sensores y la rama sana. Falta de cierre o circuito de liberación defectuoso produce alarma y condición de protección degradada, sin afirmar retención conseguida. \\\hline
F04/F05, bancos & El fuego confirmado bloquea recarga y ordena aislamiento DC mediante dispositivo apto para tensión/corriente de banco; no se usa una orden cloud como aislamiento. Se mantiene extracción de H$_2$ y evacuación. La inundación sólo se habilitará con la secuencia independiente de seguridad de gas aprobada y emisión residual comprobada; no se autoriza sellado sólo porque se haya cortado el cargador. \\\hline
Descarga / posdescarga & Presostato confirma liberación efectiva; ausencia de presión se alarma. Se mantiene restricción de acceso, salida libre y retención calculada. Reingreso por personal habilitado sólo tras ventilación posterior, medición de atmósfera y control de productos de descomposición; no se presenta FM-200 como inocuo frente a personas o superficies calientes. \\\hline
Falla / mantenimiento & Pérdida de detector, energía, tubería, compuerta, circuito o cilindro activa condición degradada. Aislamiento manual por zona requiere permiso interno, registro, aviso NOC/CLIENTE y vigilancia compensatoria; el restablecimiento sólo lo realiza el responsable local autorizado. \\\hline
\end{longtable}

Los bancos conservan ventilación segregada con la capacidad por estado de ENV16; el antecedente de cribado contempla 900 m$^3$/h efectivos por banco y 7,5978 m$^3$/h de H$_2$ como hipótesis conservadora a 60 A; GAS-28/INT-31 demuestran que ese régimen supera los umbrales propios y no autorizan esa recarga. En 36 m$^3$, sin extracción y manteniendo emisión, el incremento sería $100(7,5978/36)/60=0,35175$ puntos porcentuales/min: de 0,4 a 1,0\,\% en 1,7058 min. Esta cuenta demuestra que una retención de diez minutos y ventilación detenida no se justifican con ese escenario. Tampoco acredita mezcla uniforme ni ausencia de focos altos. El proyecto debe obtener del fabricante emisión residual después de aislar carga, ubicar medición alta por gabinete, verificar redundancia/extracción y resolver conjuntamente gas, fuego, aislamiento y retención. Mientras ese cálculo y la secuencia aprobada no estén resueltos, RT-06.17 permanece no completado para el conjunto del sitio. Se conservan la familia de agente y la sectorización elegidas; no se ofrece como cerrada una alternativa para bancos ni se certifica el diseño con esta aproximación.

\textbf{Monitor auxiliar, aviso NOC y contraparte (MON-INC-4).}\label{sec:sd4-moninc4}
Se seleccionan cuatro Schneider Electric Modicon TM221CE24R, con catorce entradas digitales de 24 V, diez salidas de relé y servidor Modbus TCP. La ficha de la variante acredita 35 VA sin expansión, alimentación 100--240 V CA y 250 mA de fuente de sensores; las 14 entradas a un máximo conservador de 7 mA requieren 98 mA por lector, dentro de ese límite. Se utilizan únicamente entradas: todas las salidas quedan sin conductores y el programa carece de lógica de actuación de incendio. No se atribuye a la certificación industrial UL 508 de estos lectores una aprobación de central de incendio \parencite[pp.~1--4]{l18-m221}.

Se asignan 48 contactos auxiliares aislados mediante relés Fike 55-048, separados de los contactos de entrada del panel y de las salidas de descarga. Se definen seis señales por zona (36): Action, Fire 1 confirmado por el panel, falla de detección, predescarga, descarga confirmada por presión y bloqueo/degradación. Las doce señales comunes son falta de alimentación A, falta B, avería de fuente base, avería suplementaria, batería baja base, batería baja suplementaria, falla agregada de lazos A (SLC1/SLC3), falla agregada de lazos B (SLC2/SLC4), mantenimiento/inhabilitación, H$_2$ alto A, H$_2$ alto B y fallo de ventilación/gas. Cada lector reserva doce entradas útiles más dos libres, total 56 entradas disponibles. Contactos de falla se cablean con lógica de circuito sano energizado; no se deriva un ``todo normal'' de la falta de energía del lector. La supervisión normativa de detección y liberación permanece en el Fike, no se sustituye por un lector industrial.

El mapa de adquisición usa F01--F06 y los seis estados en ese orden, repartidos F01/F02 en Lector1, F03/F04 en Lector2, F05/F06 en Lector3 y las doce señales comunes en Lector4. Cada programa propio publica máscara de doce bits en \texttt{\%MW100}, contador de transición en \texttt{\%MW101} y contador de vida en \texttt{\%MW102}; el integrador fija y entrega la correspondencia exacta entre esas palabras y offsets Modbus cero, con prueba de cada bit y revisión de firmware. La ACL de la VLAN de incendio admite sólo el colector local, función de lectura 03 en TCP/502; bloquea escritura/programación desde redes usuarias/WAN. Se mantienen salidas físicas desconectadas incluso si alguien altera esa ACL. La pérdida de Ethernet no cambia la central ni su autonomía.

Un colector reservado en los nodos locales existentes consulta cada lector una vez por segundo, confirma vida y conecta con la plataforma de observabilidad usando HTTPS/TLS con identidad de servicio. Se reservan 0,05 vCPU, 256 MiB de RAM y 10 GiB de disco por nodo para el servicio, con un único escritor activo y relevo mediante la autoridad local ya definida; se incorporan como recursos adicionales al inventario/capacidad. Cada transición conserva sitio, zona, señal, valor anterior/nuevo, origen, secuencia, hora UTC, tiempo monotónico, estado de sincronización y acuse. Se configura en C-LINX el enclavamiento de las seis señales de cada zona y de las doce señales comunes: cada aparición activa su contacto auxiliar durante al menos 30 s y hasta reposición local autorizada, aunque cese antes la condición física; el estado actual y el estado enclavado se presentan por separado en la central/bitácora. Fuego, predescarga, presión de descarga, H$_2$ y degradación nunca se reponen por una orden del colector. La reposición exige evidencia de persistencia local de cada estado activo y comprobación de retorno físico a normal; con WAN caída basta persistencia local y constancia del responsable, quedando el aviso remoto en cola. Si el colector está averiado, se prohíbe la reposición automática: el responsable conserva la indicación enclavada y realiza registro manual de zona/hora/causa antes de una reposición excepcional autorizada, que se reconcilia al restaurar el colector. La aceptación debe demostrar con un impulso breve de entrada que el panel produce la indicación y contacto enclavados; si el programa/relé escogido no admite esa función, se incorpora un módulo de memoria de evento compatible antes de habilitación. Es una exigencia funcional de configuración, no una prestación atribuida sin prueba al lector industrial. De este modo el sondeo de un segundo no pretende descubrir un pulso que ya desapareció. Se captura también la alarma local del Fike y se coteja con su bitácora de hasta 3.200 eventos, sin presentarla como almacenamiento ilimitado. El colector persiste antes de transmitir; conserva siete días y hasta un millón de transiciones a 1 KiB, aproximadamente 0,954 GiB, en cuota de 10 GiB con rotación, alarma al 70/85/95\,\% y prohibición de borrar eventos críticos sin acuse/archivo. Se registran estados actuales tras reinicio y se marca la ventana perdida; no se promete el detalle de todos los sensores del panel desde contactos agregados. La garantía de captura se limita a transiciones que el panel reconoce y enclava según la política anterior: un pulso inferior al tiempo de reconocimiento de sus entradas no se convierte en evento garantizado. El colector detecta cambios mientras el contacto permanece enclavado; una condición repetida durante el mismo enclavamiento sigue visible como incidente activo y sus repeticiones se cotejan con la historia del panel, sin inventar un segundo flanco.

Fuego, predescarga, descarga, H$_2$ alto y protección degradada abren incidente crítico con aviso paralelo al NOC Synaptix 24$\times$7 y a la contraparte/titular y suplente del CLIENTE inscritos en el registro de contactos. Con conectividad, la meta de entrega es 30 s desde persistencia local; recepción sin acuse a 60 s reintenta por el segundo canal, y ausencia de aceptación a dos minutos escala por teléfono al turno NOC y suplente. Las prealarmas Action mantienen identificación de zona y reciben confirmación en cinco minutos; reiteración/escalada nunca rearma ni silencia automáticamente el panel. La nómina y teléfonos se reciben antes de habilitar el servicio; no se inventan personas ni mensajes ya enviados.

Sin WAN, permanecen panel local, sirenas/estrobo por zona y anunciador cableado Fike RDU2 10-2630 en el acceso atendido de administración para la contraparte local, conectado por RS485 según ficha, sin depender de Wi-Fi ni nube. El encargado de turno ejecuta evacuación y contacto alternativo autorizado; los eventos esperan acuse en disco y se transmiten con su hora original al retornar el canal. Se genera alarma local de enlace perdido después de diez segundos y se distingue falta de comunicación de ausencia de fuego. El NOC recibe aviso remoto durante caída WAN sólo si subsiste un canal; con aislamiento exterior total no se afirma notificación instantánea. Las baterías propias aseguran el panel y anunciador; lectores/colector mantienen alimentación A/B y su autonomía se recibe separadamente \parencite[pp.~3--6]{l18-cheetah}.

Synaptix y el instalador entregan matriz causa/efecto, inventario de direcciones, esquemas de supervisión/alimentación, permisos de mantenimiento y protocolo de aceptación: humo en punto lejano de cada red, fallo de aspiración, coincidencia y alarma única, aborto/manual, simulación de disparo con actuadores sustituidos por cargas de prueba, presión simulada, compuertas, pérdida A/B, batería propia, caída WAN, disco lleno y reenvío sin duplicar incidentes. No se declara que esas pruebas, listado de conjunto, descarga o revisión técnica humana ya hayan ocurrido. Las cuatro zonas sin bancos y el monitor quedan diseñados con equipos seleccionados; la aptitud de la detección en bancos permanece pendiente; el expediente ejecutivo pendiente y los límites de EXT-4 se mantienen identificados antes de compra y puesta en servicio.
